\chapter*{Remerciements}
\addcontentsline{toc}{chapter}{Remerciements}

Je tiens a adresser mes remerciements les plus sinceres a toutes les personnes qui ont contribue, de pres ou de loin, a la realisation de ce stage et a l'aboutissement du projet SupportFlow.

Mes premiers remerciements s'adressent a l'equipe pedagogique de \Etablissement, ainsi qu'aux enseignants de la filiere \Specialite, pour la qualite de la formation dispensee et pour les competences techniques, methodologiques et humaines qui m'ont permis d'aborder ce projet avec rigueur.

J'exprime egalement ma gratitude a mon encadrant academique, \EncadrantAcademique, pour son suivi, ses conseils et son accompagnement tout au long du stage. Ses remarques ont ete precieuses pour structurer ma reflexion, clarifier la demarche du projet et renforcer la qualite du livrable final.

Je remercie aussi mon encadrant au sein de \EntrepriseAccueil, \EncadrantEntreprise, pour sa confiance, sa disponibilite et sa vision metier. Son accompagnement m'a aide a mieux comprendre les enjeux concrets d'une plateforme de support, en particulier les questions de priorisation, de workflow, de supervision et de qualite de service.

Je souhaite egalement remercier l'ensemble des collaborateurs, collegues et personnes ressources qui ont facilite mon integration et enrichi ce stage par leurs retours, leurs suggestions et leurs echanges. Les discussions autour des besoins utilisateurs, des choix techniques, de la securite, du BPMN, de la GED, de l'IA et du deploiement ont fortement contribue a la maturation de la solution.

Enfin, je remercie ma famille et mes proches pour leur soutien moral, leur patience et leurs encouragements tout au long de cette periode exigeante.

Ce rapport constitue l'aboutissement d'un travail d'analyse, de conception, de realisation, de correction et de validation. Il reflete non seulement le produit realise, mais aussi l'experience humaine et professionnelle acquise durant ce stage.

